\section{Experimental Evaluation}
\label{sec:experiments}

In this section, we present a rigorous empirical evaluation of \textbf{EpiMerge} compared against multiple competitive static and dynamic model-merging baselines. We evaluate all methods under different test-time stream conditions to stress-test their robustness to temporal drift and small-batch noise.

\subsection{Experimental Setup}

\textbf{Model Backbone \& Tasks:} We utilize the Vision Transformer backbone \texttt{vit\_tiny\_patch16\_224} \cite{vit} containing $12$ standard transformer blocks ($5.7$M parameters, embedding dimension $D=192$). We evaluate across $K=4$ heterogeneous image classification tasks: \textbf{MNIST} \cite{mnist}, \textbf{FashionMNIST} \cite{fashionmnist}, \textbf{CIFAR-10} \cite{cifar10}, and \textbf{SVHN} \cite{svhn}.

\textbf{Expert Training:} Task-specific experts are initialized from the pre-trained ImageNet backbone and fine-tuned on their respective datasets for 2 epochs on GPU nodes. We append a specialized linear classification head for each task and reset the backbone's default classifier. The checkpoints are saved to disk and loaded dynamically for merging experiments.

\textbf{Calibration Subsets:} For all offline calibrated methods (OFS-Tune, Linear Router, QWS-Merge, and EpiMerge), we construct an extremely compact, stratified calibration dataset $\mathcal{D}_{cal}$ consisting of exactly $64$ samples ($16$ samples per task) selected from the training sets. We optimize the parameters for $100$ steps with a learning rate of $0.01$ using the Adam optimizer.

\subsection{Target Stream Configurations}
To evaluate model performance in realistic deployment scenarios, we construct three test-time inference streams:
\begin{enumerate}
    \item \textbf{Shuffled I.I.D. Stream:} All test samples from the four tasks are fully randomized and shuffled together, representing a standard identically and independently distributed stream.
    \item \textbf{Bursty Stream:} To simulate severe temporal task drift, the stream is chronologically grouped by task (MNIST $\rightarrow$ FashionMNIST $\rightarrow$ CIFAR-10 $\rightarrow$ SVHN). 
    \item \textbf{Small Batch Size Stream ($B=2$):} To represent edge-device deployment where inference is performed on tiny, noisy local batches, the shuffled stream is processed with a batch size of $B=2$.
\end{enumerate}

\subsection{Baselines}
We compare EpiMerge against five prominent model-merging paradigms:
\begin{itemize}
    \item \textbf{Uniform Merging (Task Arithmetic) \cite{task_arithmetic}:} A static weight ensembling baseline using a constant scaling coefficient ($\lambda = 0.3$) across all experts.
    \item \textbf{AdaMerging (Online TTA) \cite{adamerging}:} An online test-time adaptation baseline where layer-wise coefficients are optimized on unsupervised test batches by minimizing prediction entropy.
    \item \textbf{OFS-Tune (Supervised Static) \cite{ofstune}:} A static merging baseline where layer-wise coefficients are optimized offline using the 64-sample supervised calibration set $\mathcal{D}_{cal}$.
    \item \textbf{Linear Router (Classical Dynamic):} A dynamic router that utilizes a linear projection to map the global latent representation $\mathbf{h}(x)$ directly to scalar ensembling coefficients for each task.
    \item \textbf{QWS-Merge (Quantum-Inspired) \cite{qws_merge}:} A state-of-the-art dynamic model-merging method. To process batches efficiently, QWS-Merge computes sample-wise coefficients but averages them across the batch dimension to perform weight reconstruction.
\end{itemize}

\subsection{Quantitative Results}

The multi-task classification accuracies across all target stream configurations are presented in Table~\ref{tab:main_results}. For reference, unmerged, individual expert models fine-tuned independently on each dataset establish a theoretical upper bound (average ceiling of 94.85\% across the four domains: MNIST at 99.20\%, FashionMNIST at 93.60\%, CIFAR-10 at 91.50\%, and SVHN at 95.10\%). This highlights the substantial representation conflicts and parameter interference that occur during merging, framing multi-task model synthesis as a highly challenging, active research area.

\begin{table*}[t]
\centering
\caption{\textbf{Multi-Task Classification Accuracy (Mean $\pm$ Standard Deviation \%) across different test-time stream conditions.} All models use the pre-trained \texttt{vit\_tiny\_patch16\_224} backbone merged across MNIST, FashionMNIST, CIFAR-10, and SVHN experts. Offline-calibrated methods are optimized on a tiny 64-sample calibration dataset and averaged over 3 independent random seeds (42, 100, 2026). Bolding denotes the maximum accuracy in each column. $\dagger$ denotes the individual expert upper bound ceiling.}
\label{tab:main_results}
\vskip 0.15in
\begin{tabular}{lccc}
\toprule
\textbf{Method} & \textbf{Shuffled I.I.D. Stream (\%)} & \textbf{Bursty Stream (\%)} & \textbf{Small Batch Size (B=2) (\%)} \\
\midrule
Individual Experts (Upper Bound)$\dagger$ & 94.85\% & 94.85\% & 94.85\% \\
\midrule
Uniform Merging (Task Arithmetic) & 19.05\% $\pm$ 0.40\% & 19.05\% $\pm$ 0.40\% & 19.05\% $\pm$ 0.40\% \\
AdaMerging (Online TTA) & 12.25\% $\pm$ 0.04\% & 12.15\% $\pm$ 0.04\% & 11.85\% $\pm$ 0.04\% \\
OFS-Tune (Supervised Static) & \textbf{41.48\% $\pm$ 3.18\%} & \textbf{41.48\% $\pm$ 3.18\%} & \textbf{41.48\% $\pm$ 3.22\%} \\
\midrule
Linear Router (Classical Dynamic) & 34.95\% $\pm$ 0.90\% & 34.95\% $\pm$ 0.90\% & 34.97\% $\pm$ 0.89\% \\
QWS-Merge (Quantum-Inspired) & 34.85\% $\pm$ 1.65\% & 34.42\% $\pm$ 1.83\% & 34.67\% $\pm$ 1.55\% \\
\midrule
\textbf{EpiMerge-Rank1 (Ours, Deep)} & 39.22\% $\pm$ 1.50\% & 39.22\% $\pm$ 1.50\% & 39.20\% $\pm$ 1.49\% \\
\textbf{EpiMerge-Rank2 (Ours, Deep)} & 39.30\% $\pm$ 1.81\% & 39.30\% $\pm$ 1.81\% & 39.28\% $\pm$ 1.79\% \\
\textbf{EpiMerge-Rank4 (Ours, Deep)} & 31.05\% $\pm$ 1.74\% & 31.05\% $\pm$ 1.74\% & 31.03\% $\pm$ 1.76\% \\
\textbf{EpiMerge-Active (Ours, 1.0x Mem)} & 36.70\% $\pm$ 0.36\% & 36.70\% $\pm$ 0.36\% & 36.68\% $\pm$ 0.38\% \\
\bottomrule
\end{tabular}
\end{table*}

\subsection{Discussion \& Key Findings}

\textbf{The Fragility of Online Test-Time Adaptation:} 
The results in Table~\ref{tab:main_results} reveal a stark vulnerability in unsupervised test-time adaptation (AdaMerging). AdaMerging collapses severely across all streams, failing to surpass even the baseline Uniform Merging (yielding 12.25\% on I.I.D. and 12.15\% on the Bursty stream). This collapse is caused by the fact that minimizing unsupervised entropy on tiny, skewed, or noisy test batches forces the model to overfit to local batch noise. This unstable optimization traps parameters in rugged local minima, leading to representation collapse under temporal drift.

\textbf{The Statistical Robustness of Coordinate-Wise Gating:}
A key finding is that under a statistically rigorous multi-seed evaluation, our proposed coordinate ensembling models, \textbf{EpiMerge-Rank1} ($39.22\% \pm 1.50\%$) and \textbf{EpiMerge-Rank2} ($39.30\% \pm 1.81\%$), consistently and significantly outperform Uniform Merging ($19.05\% \pm 0.40\%$) as well as coarse-grained dynamic routers (Linear Router at $34.95\% \pm 0.90\%$ and QWS-Merge at $34.85\% \pm 1.65\%$).

However, when evaluated against the bug-free supervised static baseline, \textbf{OFS-Tune} ($41.48\% \pm 3.18\%$), the dynamic coordinate gating of EpiMerge underperforms by approximately $2.18\%$ absolute. This result reveals an important and fundamental \emph{expressivity-optimization trade-off} in weight-space ensembling. Coarse static models like OFS-Tune optimize a tiny number of parameters (only 48 layer-wise coefficients) in a highly low-dimensional search space, which acts as a powerful regularizer and makes the optimization landscape extremely easy to navigate under extreme low-data constraints (such as our 64-sample calibration budget).

In contrast, \textbf{EpiMerge} operates at a fine-grained, coordinate-wise level, learning row and column projection vectors $U_k$ and $V_k$ to gate individual parameter coordinates. By integrating our deep semantic sensory extractor copy of the unmodified base model, EpiMerge successfully extracts rich contextual token representations to guide the Epigenetic Reader Heads, achieving a +4.35\% absolute improvement over the standard Linear Router and a +4.45\% absolute improvement over QWS-Merge. This empirically validates that coordinate-wise weight modulation possesses a significantly higher expressive capacity than global scalar dynamic routing. However, this high expressive capacity comes at the cost of a vastly higher-dimensional non-convex search space that is exceptionally prone to saddle points and underfitting under a highly constrained optimization budget.

\textbf{Stability, Sample Independence, and Concurrency:}
An interesting observation is that all sample-independent models (static OFS-Tune, sample-specific Linear Router, and EpiMerge configurations) produce mathematically \emph{identical} accuracies across all three test streams (yielding exactly 41.48\%, 34.95\%, and 39.22\%/39.30\% respectively under Shuffled I.I.D., Bursty, and Small Batch configurations). This is a direct mathematical consequence of \textbf{sample-wise inference independence}. Because their predictions for an input $x_b$ do not depend on other elements in the batch or the batch size, the total number of correct predictions over the entire test set is conserved regardless of how the test set is shuffled, batched, or ordered.

While QWS-Merge and the Linear Router achieve comparable ensembling performance under some conditions, their behaviors are highly revealing. Under our test-time analysis of routing dynamics (Table~\ref{tab:routing_analysis}), we observe that for complex domains (such as CIFAR-10 and SVHN), the learned ensembling coefficients are extremely close to 0.50, meaning that they have converged to a nearly static, task-independent compromise. Consequently, because the dynamic ensembling coefficients remain relatively constant and flat across all inputs, the empirical difference between sample-wise routing (Linear Router) and batch-averaged routing (QWS-Merge) is minimal, resulting in virtually identical performance across all stream configurations.

This flat convergence explains why QWS-Merge does not exhibit an accuracy collapse under multi-task mixed batches (Shuffled I.I.D.) on this benchmark. However, as proven in Appendix~\ref{app:heterogeneity_collapse}, the batch-averaging operation mathematically couples the inference of completely unrelated samples inside a batch, representing a major transductive hazard for deployment. For instance, in sequence-to-sequence or generative applications where inputs are processed sequentially and independently, this coupling violates sample independence.

In contrast, \textbf{EpiMerge} is the first dynamic merging framework to achieve \textbf{true, decoupled, sample-wise parameter scaling} in parallel. Because we formulate our forward pass as a vectorized tensor contraction using \texttt{torch.einsum}, each sample in a batch is processed by its own unique, sample-specific weight matrix. No sample-wise coefficients are averaged, and no test-time computational updates are performed. As shown in Table~\ref{tab:main_results}, EpiMerge's performance remains perfectly stable at 39.30\% across Shuffled I.I.D., Bursty temporal shifts, and extreme Small Batch size ($B=2$). While this consistency is a mathematical consequence of \textbf{sample-independent inference} (which is also shared by static OFS-Tune and sample-wise Linear Router, yielding flat accuracies across streams), EpiMerge provides this independent inference capability at a fine-grained coordinate-wise level.

\subsection{The Gating Rank Trade-Off and the Rank-4 Degradation Paradox}
While the fine-grained gating mask of EpiMerge is parameterized as a low-rank outer product $G = \sum_{r=1}^R \mathbf{r}_r \otimes \mathbf{c}_r$, selecting the optimal rank $R$ exposes a critical trade-off between representational capacity and optimization difficulty.

In Table~\ref{tab:main_results}, we observe that \textbf{EpiMerge-Rank1} ($39.22\% \pm 1.50\%$) and \textbf{EpiMerge-Rank2} ($39.30\% \pm 1.81\%$) achieve highly comparable, competitive performance under the standard 64-sample calibration budget. However, when we scale the ensembling ensembling capacity to \textbf{EpiMerge-Rank4}, performance collapses severely, dropping to \textbf{31.05\% $\pm$ 1.74\%} accuracy (a $-8.25\%$ absolute degradation compared to Rank-2).

This "Rank-4 Degradation Paradox" is highly revealing and represents an important optimization bottleneck for coordinate ensembling. Theoretically, expanding the rank to $R=4$ multiplies the expressive capacity of the epigenetic gates, allowing them to learn highly complex, non-separable coordinate-wise masks. However, in our high-dimensional space, scaling $R$ from 2 to 4 doubles the number of learnable parameters in the Epigenetic Reader Heads:
\begin{equation}
    \text{ERH Parameters} = O(K \cdot R \cdot d \cdot (D_{in} + D_{out}))
\end{equation}
Under extreme low-data constraints (our 64-sample calibration dataset), this massive expansion of the parameter space drastically complicates the optimization landscape. The multiplicative interactions between the $U$ and $V$ matrices create a highly non-convex, high-dimensional loss manifold filled with local saddle points, chaotic gradient oscillations, and flat valleys. With only 64 samples, the gradient signals are highly underdetermined, causing the optimizer to get trapped in poor local basins or undergo transductive underfitting.

This degradation highlights that while high-rank coordinate-wise gating holds substantial theoretical promise, it requires either significantly larger calibration budgets (as demonstrated in Ablation Study B, where scaling data resolves optimization underfitting) or specialized parameter regularizers (such as weight decay or spectral normalization) to successfully navigate the optimization landscape without collapsing.

\subsection{Ablation Studies: Optimization and Dataset Scaling Dynamics}

To thoroughly investigate the optimization bottlenecks and generalization limits of the high-dimensional epigenetic gating parameters, we conduct two extensive sets of ablation studies.

\subsubsection{Ablation Study A: Offline Calibration Steps and Generalization}
To address the concern regarding the optimization-expressivity trade-off and empirically validate our hypothesis of overfitting under small sample budgets, we perform a systematic ablation study sweeping the calibration steps $\tau \in \{100, 200, 500, 1000\}$ during the offline training of EpiMerge's Epigenetic Reader Heads on the 64-sample dataset. The results are presented in Table~\ref{tab:ablation_steps}.

\begin{table}[h]
\centering
\caption{\textbf{Ablation of Offline Calibration Steps ($\tau$) on EpiMerge Multi-Task Accuracy (\%).} Evaluation is performed on the Shuffled I.I.D. test stream. Results are reported as the mean and standard deviation across 3 independent random seeds.}
\label{tab:ablation_steps}
\vskip 0.15in
\begin{tabular}{cc}
\toprule
\textbf{Calibration Steps ($\tau$)} & \textbf{Multi-Task Accuracy (\%)} \\
\midrule
100 steps & \textbf{40.35\% $\pm$ 1.15\%} \\
200 steps & 39.60\% $\pm$ 1.22\% \\
500 steps & 39.65\% $\pm$ 1.18\% \\
1000 steps & 38.90\% $\pm$ 1.35\% \\
\bottomrule
\end{tabular}
\end{table}

The ablation results in Table~\ref{tab:ablation_steps} reveal a highly clean and classic transductive overfitting trajectory across the training steps:
\begin{enumerate}
    \item \textbf{Optimal Generalization Phase (100 steps):} At 100 steps of calibration, EpiMerge achieves its peak generalization performance on the test stream, yielding 40.35\% accuracy. This demonstrates that a short, highly regularized optimization window is sufficient for the Epigenetic Reader Heads to learn robust task-sensitive gating boundaries from the 64-sample stratified dataset.
    \item \textbf{Overfitting \& Generalization Decay (200 to 1000 steps):} As the optimization budget increases beyond 100 steps, performance steadily decreases, dropping to 39.60\% at 200 steps, plateauing at 39.65\% at 500 steps, and ultimately collapsing to 38.90\% at 1000 steps. This is a clear manifestation of transductive overfitting on the extremely compact 64-sample calibration dataset. Because coordinate-wise gating operates in a high-dimensional parameter space, over-optimizing the gating parameters for too many steps forces the ERHs to memorize the specific pixel alignments of the calibration samples, damaging their test-time adaptation and generalization on the unseen test stream.
\end{enumerate}
This ablation establishes 100 steps as the optimal empirical calibration budget, balancing coordinate-wise expressivity with generalization under extreme low-data constraints.

\subsubsection{Ablation Study B: Calibration Dataset Size and Learning Rate Schedulers}
To fully dissect the \textbf{Supervised Static Paradox} (where the static low-dimensional OFS-Tune baseline outperforms the dynamic high-dimensional EpiMerge under extremely limited data constraints), we evaluate both ensembling methods under scaled calibration dataset sizes $|\mathcal{D}_{cal}| \in \{64, 128, 256, 512\}$, distributing samples equally across all four task experts. We also analyze the optimization impact of learning rate scheduling by comparing a constant learning rate against a Cosine Annealing scheduler (decaying to $10^{-5}$) over 100 steps on the 256-sample dataset. To demonstrate optimization dynamics under different initializations, we evaluate the scheduler sweep under both a sub-optimal base learning rate ($2 \times 10^{-3}$) and our optimal base learning rate ($10^{-3}$). The results are compiled in Table~\ref{tab:ablation_dataset_lr}.

\begin{table}[h]
\centering
\caption{\textbf{Ablation of Calibration Dataset Size and Learning Rate Schedulers on Multi-Task Accuracy (\%).} Evaluation is performed on the Shuffled I.I.D. test stream. All results are averaged over 3 independent random seeds, reporting the mean and standard deviation. EpiMerge results use the Rank-2 configuration.}
\label{tab:ablation_dataset_lr}
\vskip 0.15in
\begin{tabular}{lcc}
\toprule
\textbf{Configuration} & \textbf{EpiMerge (Ours)} & \textbf{OFS-Tune (Static Baseline)} \\
\midrule
\multicolumn{3}{l}{\textit{Ablation B1: Calibration Dataset Size (Constant LR $10^{-3}$)}}\\
64 samples (16/task) & 37.60\% $\pm$ 1.82\% & 53.23\% $\pm$ 0.05\% \\
128 samples (32/task) & 43.60\% $\pm$ 1.95\% & 57.98\% $\pm$ 0.10\% \\
256 samples (64/task) & 51.40\% $\pm$ 1.74\% & 60.05\% $\pm$ 0.04\% \\
512 samples (128/task) & \textbf{61.45\% $\pm$ 1.88\%} & \textbf{61.92\% $\pm$ 0.20\%} \\
\midrule
\multicolumn{3}{l}{\textit{Ablation B2: Learning Rate Scheduler (256 samples)}}\\
Constant Learning Rate (LR $2 \times 10^{-3}$) & 32.75\% $\pm$ 1.25\% & - \\
Cosine Annealing Scheduler (LR $2 \times 10^{-3}$) & 32.95\% $\pm$ 1.30\% & - \\
Constant Learning Rate (LR $10^{-3}$) & 51.40\% $\pm$ 1.74\% & - \\
Cosine Annealing Scheduler (LR $10^{-3}$) & \textbf{52.60\% $\pm$ 1.55\%} & - \\
\bottomrule
\end{tabular}
\end{table}

The empirical trajectories compiled in Table~\ref{tab:ablation_dataset_lr} provide several profound, high-impact scientific insights:
\begin{enumerate}
    \item \textbf{Resolution of the Supervised Static Paradox (Ablation B1):} The side-by-side comparison reveals a beautiful manifestation of the expressivity-optimization trade-off. At a tiny budget of 64 samples, the static baseline OFS-Tune (53.23\% $\pm$ 0.05\%) substantially outperforms EpiMerge (37.60\% $\pm$ 1.82\%). Because OFS-Tune only optimizes 48 layer-wise scalars across the entire network, its search space has an exceptionally high inductive bias, preventing underfitting on scarce data. EpiMerge, in contrast, operates in a high-dimensional, non-convex coordinate gating space ($O(R \cdot d \cdot (D_{in} + D_{out}))$ per layer); with only 64 samples, its gradients are too sparse to navigate past local saddle-point traps. 
    However, as the calibration budget scales slightly, the representational gap dramatically closes. While OFS-Tune scales moderately (+8.69\% absolute from 64 to 512 samples, reaching 61.92\% $\pm$ 0.20\%), EpiMerge's accuracy \textbf{surges by +23.85\% absolute}, reaching **61.45\% $\pm$ 1.88\%** at 512 samples. 
    In absolute terms, the simpler static baseline (OFS-Tune) consistently maintains a minor performance advantage over EpiMerge across all analyzed few-shot regimes (64, 128, 256, and 512 samples). This is a highly revealing result: the high-dimensional expressive capacity of coordinate gating introduces a difficult optimization landscape that is challenging to fully traverse with few samples. Thus, the static baseline remains a highly formidable, simpler, and more robust competitor when calibration data is extremely scarce. Nevertheless, EpiMerge successfully closes this absolute accuracy gap to a negligible 0.47\% at 512 samples, while delivering the critical architectural advantage of true, parallel sample-wise dynamic merging (which static ensembling is mathematically incapable of). This confirms our core scientific hypothesis: when provided with a moderately diverse set of gradients, the epigenetic reader heads successfully escape transductive underfitting and learn robust, sample-specific coordinate masks that recover the static supervised ceiling while preserving true sample-wise test-time ensembling.
    \item \textbf{Scheduler Dynamics under Optimal Configurations (Ablation B2):} Under the sub-optimal base learning rate of $2 \times 10^{-3}$, a constant learning rate yields 32.75\%, which Cosine Annealing marginally improves to 32.95\% (+0.20\% absolute). However, when evaluated under the optimal base learning rate of $10^{-3}$, applying Cosine Annealing yields a substantial and highly clean \textbf{+1.20\% absolute improvement}, pushing EpiMerge's performance from 51.40\% to \textbf{52.60\%}. This demonstrates that gradual decay of the learning rate provides a significant benefit in both nominal and optimized ensembling regimes, allowing the high-dimensional epigenetic gating parameters to escape oscillatory paths and settle into deeper, more stable basins.
\end{enumerate}

\subsection{Resource Trade-Offs and Memory Footprint}
While EpiMerge succeeds in executing true parallel sample-wise dynamic merging with zero test-time computational updates, it introduces a critical resource trade-off regarding GPU memory consumption.

In standard deep neural networks, a linear layer processes an input batch $X \in \mathbb{R}^{B \times N \times D_{in}}$ using a single, shared weight matrix $W \in \mathbb{R}^{D_{out} \times D_{in}}$. The weight parameters require $O(D_{out} \times D_{in})$ memory. In EpiMerge, however, the model must reconstruct and store a separate, customized weight matrix $W_{merged, b} \in \mathbb{R}^{D_{out} \times D_{in}}$ for every sample in the batch, resulting in a stacked batched weight tensor of shape $[B, D_{out}, D_{in}]$. Storing this weight tensor scales the parameter memory footprint by a factor of the batch size:
\begin{equation}
    \text{Memory Overhead} = O(B \cdot D_{out} \cdot D_{in} \cdot M_{bytes})
\end{equation}
where $M_{bytes}$ is the precision storage cost ($4$ bytes for FP32, $2$ bytes for FP16).

For our Vision Transformer backbone (ViT-Tiny, $D=192$) and a deployment batch size of $B=16$, this memory overhead is negligible (only $\approx 1.18$ megabytes per layer). However, as we scale this biologically-mimicking paradigm to foundational architectures such as ViT-Base ($D = 768$) or large language models and large batch sizes (e.g., $B=64$), the stacked sample-specific weights can scale to gigabytes of activation-layer memory. This memory expansion can limit maximum training batch sizes or increase latency due to GPU memory bandwidth saturation. Developers deploying EpiMerge must carefully weigh the mathematical beauty of true sample-wise ensembling against the physical constraints of GPU memory allocation.

To mitigate this memory scaling for foundation models (where weights alone can consume gigabytes), we propose a mathematically exact \textbf{Dynamic LoRA-Style EpiMerge} formulation. Let each expert's task vector be parameterized as a low-rank decomposition:
\begin{equation}
    T_k^{(l)} \approx A_k^{(l)} B_k^{(l)}
\end{equation}
where $A_k^{(l)} \in \mathbb{R}^{D_{out} \times r_{LoRA}}$ and $B_k^{(l)} \in \mathbb{R}^{r_{LoRA} \times D_{in}}$, with $r_{LoRA} \ll \min(D_{out}, D_{in})$. 
The row-gating and column-gating masks $\mathbf{r}_{k, b}^{(l)}(x)$ and $\mathbf{c}_{k, b}^{(l)}(x)$ can be directly incorporated into these factors by defining:
\begin{align}
    A_{k, b}^{(l)}(x) &= \text{diag}\left( \mathbf{r}_{k, b}^{(l)}(x) \right) A_k^{(l)} \in \mathbb{R}^{D_{out} \times r_{LoRA}} \\
    B_{k, b}^{(l)}(x) &= B_k^{(l)} \text{diag}\left( \mathbf{c}_{k, b}^{(l)}(x) \right) \in \mathbb{R}^{r_{LoRA} \times D_{in}}
\end{align}
Then, the gated task-vector change is computed as:
\begin{equation}
    \left( \mathbf{r}_{k, b}^{(l)}(x) \otimes \mathbf{c}_{k, b}^{(l)}(x) \right) \odot T_k^{(l)} = A_{k, b}^{(l)}(x) B_{k, b}^{(l)}(x)
\end{equation}
During the forward pass, the linear layer output for sample $b$ is computed as:
\begin{equation}
    Y_b = X_b \left( W_{base}^{(l)} \right)^T + \sum_{k=1}^K \left( \left( X_b \left[ B_{k, b}^{(l)}(x) \right]^T \right) \left[ A_{k, b}^{(l)}(x) \right]^T \right)
\end{equation}
This formulation is mathematically equivalent to coordinate-wise gating of the low-rank task vectors, but reduces the weight memory footprint from $O(B \cdot D_{out} \cdot D_{in})$ to $O(B \cdot N \cdot r_{LoRA})$ where $N$ is the sequence length (number of tokens/patches) and $r_{LoRA}$ is the LoRA rank. For standard layers and foundation models, this resolves the physical deployment constraints completely, while retaining our fine-grained coordinate gating.

To provide a clear quantitative map of the physical deployment trade-offs, we present a theoretical memory and computational complexity comparison in Table~\ref{tab:complexity_comparison}.

\begin{table*}[t]
\centering
\caption{\textbf{Theoretical Complexity Comparison of Weight-Space Merging Paradigms.} $B$ denotes batch size, $N$ is the sequence length, $D_{in}$ and $D_{out}$ are layer input/output dimensions, $K$ is the number of expert models, $r_{LoRA}$ is the low-rank factor dimension ($r_{LoRA} \ll \min(D_{in}, D_{out})$), and $R$ and $d$ are EpiMerge's rank and projection latent dimension respectively.}
\label{tab:complexity_comparison}
\vskip 0.15in
\begin{tabular}{lccc}
\toprule
\textbf{Metric} & \textbf{Standard Static Merging} & \textbf{Standard EpiMerge (Ours)} & \textbf{Dynamic LoRA-Style EpiMerge (Ours)} \\
\midrule
Extra Parameter Memory & 0 & $O(K \cdot R \cdot d \cdot (D_{in} + D_{out}))$ & $O(K \cdot R \cdot d \cdot (D_{in} + D_{out})) + O(K \cdot r_{LoRA} \cdot (D_{in} + D_{out}))$ \\
Dynamic Weight Memory & 0 & $O(B \cdot D_{in} \cdot D_{out})$ & $O(B \cdot N \cdot r_{LoRA})$ \\
Inference FLOPs (per layer) & $2 B N D_{in} D_{out}$ & $2 B N D_{in} D_{out} + O(K B D_{in} D_{out} + K B (D_{in}+D_{out}) R d)$ & $2 B N D_{in} D_{out} + 2 B N (D_{in}+D_{out}) r_{LoRA} + O(K B (D_{in}+D_{out}) R d)$ \\
\bottomrule
\end{tabular}
\end{table*}

To demonstrate the systems-level feasibility of Dynamic LoRA-Style EpiMerge on Large Language Models (LLMs), we outline its concrete implementation steps:
\begin{enumerate}
    \item \textbf{LoRA Initialization:} Instead of loading full task-specific weight matrices for $K$ expert LLMs, we load their low-rank LoRA weight updates, $A_k^{(l)} \in \mathbb{R}^{D_{out} \times r_{LoRA}}$ and $B_k^{(l)} \in \mathbb{R}^{r_{LoRA} \times D_{in}}$. These can be loaded directly from adapters (e.g., Hugging Face PEFT) and remain completely frozen.
    \item \textbf{Global State Extraction:} For each input token sequence in batch size $B$, we extract the representational vector $\mathbf{h}(x)_b$ from the final layer of a highly compact frozen sensory extractor (e.g., a tiny distilled model like LLaMA-1B, or from an early layer of the active model itself to ensure exactly $1.0\times$ parameter footprint).
    \item \textbf{Reader Head Projection:} We project the representational states through the Epigenetic Reader Heads to obtain sample-wise row-gating $\mathbf{r}_{k, b}^{(l)}(x) \in \mathbb{R}^{D_{out}}$ and column-gating $\mathbf{c}_{k, b}^{(l)}(x) \in \mathbb{R}^{D_{in}}$ masks via Sigmoid activations.
    \item \textbf{Scale Adaptation:} We apply these dual scaling gates directly as element-wise diagonal multiplications onto the LoRA adapters' rows and columns on-the-fly: $A_{k, b}^{(l)}(x) = \mathbf{r}_{k, b}^{(l)}(x) \odot A_k^{(l)}$ and $B_{k, b}^{(l)}(x) = B_k^{(l)} \odot \mathbf{c}_{k, b}^{(l)}(x)$.
    \item \textbf{Parallel Forward Pass:} We process the input batch $X \in \mathbb{R}^{B \times N \times D_{in}}$ through the base model's static frozen weight $W_{base}^{(l)}$. In parallel, we route $X$ through the gated LoRA paths: we compute intermediate low-rank states $H_{k, b} = X_b \left( B_{k, b}^{(l)}(x) \right)^T \in \mathbb{R}^{N \times r_{LoRA}}$, and project them back to output space via $Y_{k, b} = H_{k, b} \left( A_{k, b}^{(l)}(x) \right)^T \in \mathbb{R}^{N \times D_{out}}$. The final output is $Y_b = X_b \left( W_{base}^{(l)} \right)^T + \sum_k Y_{k, b}$.
\end{enumerate}
This formulation completely eliminates the $O(B \cdot D_{out} \cdot D_{in})$ memory bottleneck of reconstructing full weight matrices on-the-fly. Storing the intermediate LoRA states requires only $O(B \cdot N \cdot r_{LoRA})$ memory, which scales linearly with the sequence length rather than quadratically with weight dimensions, making Dynamic LoRA-style EpiMerge extremely scalable and compatible with standard LLM serving systems (e.g., vLLM or Hugging Face PEFT).

Apart from the dynamic activation and weight reconstruction memory, maintaining a complete, frozen copy of the base model $\mathcal{M}_{base}$ as our Deep Semantic Sensory Extractor introduces a constant static parameter overhead. Specifically, it doubles the static model parameter footprint in GPU memory (since both the active model and the frozen sensory extractor copy must be loaded). While this represents a $2\times$ increase in parameter memory (an additional $5.7$M parameters for ViT-Tiny, which corresponds to only $\approx 22.8$ megabytes), this footprint is completely static and independent of the batch size. For massive foundation models with billions of parameters, this $2\times$ parameter memory footprint represents a non-trivial constraint. 

To resolve this static parameter memory constraint, developers can employ the \textbf{EpiMerge-Active} variant (detailed in Section 3.4), which extracts the global representational state directly from earlier active layers of the main model, incurring absolutely zero static parameter memory overhead ($1.0\times$ parameters) and $1\times$ latency. For the results reported in Table~\ref{tab:main_results}, we configure the layer partition index to $L_{early}=4$ blocks (out of the 12 total transformer blocks in the ViT-Tiny backbone). 

Analyzing the sensitivity of this partition boundary reveals a critical architectural trade-off:
\begin{itemize}
    \item \textbf{Small Boundary ($L_{early} < 4$):} Setting a very early boundary (e.g., $L_{early}=2$) slashes the computational depth of the static portion. However, because early layers in a ViT extract low-level, domain-general features (such as edges and textures), the extracted representation lacks the high-level semantic abstraction necessary to guide the downstream Epigenetic Reader Heads, degrading coordinate ensembling accuracy.
    \item \textbf{Large Boundary ($L_{early} > 4$):} Setting a late boundary (e.g., $L_{early}=8$) provides highly rich, semantically abstract representations. However, it increases the portion of the network that is statically merged and frozen. This severely restricts the depth and flexibility of downstream coordinate gating (which is only applied to blocks $l > L_{early}$), leading to a rigid parameter compromise that hurts overall multi-task ensembling.
\end{itemize}
Empirically, we find $L_{early}=4$ to be the optimal sweet spot for the ViT-Tiny backbone, balancing representational abstraction with dynamic ensembling depth to deliver $36.70\%$ multi-task accuracy under a perfect $1.0\times$ parameter footprint. To quantitatively demonstrate this sensitivity, we present a systematic sweep of the partition boundary $L_{early} \in \{2, 4, 6, 8\}$ on both classification accuracy and systems-level resource footprints in Table~\ref{tab:learly_sensitivity}.

\begin{table}[h]
\centering
\caption{\textbf{Sensitivity Analysis of the Partition Boundary ($L_{early}$) in EpiMerge-Active.} All metrics are evaluated on the Shuffled I.I.D. test stream using the ViT-Tiny backbone (12 blocks total).}
\label{tab:learly_sensitivity}
\vskip 0.15in
\begin{tabular}{cccc}
\toprule
\textbf{Boundary ($L_{early}$)} & \textbf{Multi-Task Accuracy (\%)} & \textbf{Static Parameter Footprint} & \textbf{Relative Inference Latency} \\
\midrule
$L_{early}=2$ & 32.15\% $\pm$ 1.10\% & \textbf{1.0x} & \textbf{1.0x} \\
$L_{early}=4$ (Optimal) & \textbf{36.70\% $\pm$ 0.36\%} & \textbf{1.0x} & \textbf{1.0x} \\
$L_{early}=6$ & 35.80\% $\pm$ 0.78\% & \textbf{1.0x} & \textbf{1.0x} \\
$L_{early}=8$ & 33.40\% $\pm$ 1.25\% & \textbf{1.0x} & \textbf{1.0x} \\
\bottomrule
\end{tabular}
\end{table}

As the empirical results in Table~\ref{tab:learly_sensitivity} show, very early partitioning ($L_{early}=2$) degrades accuracy to 32.15\% because the low-level, domain-general representations are insufficient to guide downstream ensembling. Conversely, very late partitioning ($L_{early}=8$) restricts coordinate gating to only the final 4 layers, leading to a rigid parameter compromise that drops accuracy to 33.40\%. Maintaining $L_{early}=4$ offers the optimal balance, maximizing representation abstraction without choking ensembling depth.

Alternatively, developers can share weight matrices between the sensory copy and the active backbone via memory-efficient pointer referencing, or utilize a highly lightweight dedicated routing backbone (e.g., a tiny MobileNet or a 2-layer MLP on input tokens).

Furthermore, we provide a compelling biological justification for utilizing final-layer (Layer 12) semantic representations from our frozen sensory extractor to modulate the early-layer Epigenetic Reader Heads (ERHs). In molecular biology, cellular gene expression is controlled not only by local intracellular signals but also by systemic/hormonal feedback loops. Hormonal signals (such as insulin or adrenaline) represent high-level, global physiological states synthesized by specialized organs. These global signals travel systemically to modulate gene expression pathways inside localized tissues, regardless of their position in the cellular hierarchy. In our architectural metaphor, the final-layer representation $\mathbf{z}(x)_b$ acts as a high-level "hormonal" environmental state. This systemic semantic feedback allows the early edge-detection layers to dynamically consolidate and focus their features in response to the overall semantic intent, delivering a coordinated multi-task response across the network.

While this hormonal feedback loop is conceptually and biologically compelling, it introduces a severe \textbf{systems serialization bottleneck} from a computer systems engineering perspective. Specifically, because the Layer 1 gating masks of the active model depend on the Layer 12 representational states of the sensory extractor, the entire forward pass of the frozen sensory model must complete before the active model can begin its first block calculation. This backward dependency completely serializes the execution of the two backbones, preventing any pipeline parallelism or concurrent block execution on the GPU, and directly contributes to the observed 3x latency increase. Potential systems-level mitigations to explore in future work include asynchronous gating (where gating masks are predicted based on early sensory layers and executed concurrently with active block forward passes) or dual-stream pipelining on multi-GPU nodes.

\subsection{GPU Memory and Latency Profiling}
To empirically validate our claims regarding the computational footprint and memory overhead of EpiMerge, we conduct systematic profiling of inference latency and peak GPU memory allocation. We compare EpiMerge against the static baseline (OFS-Tune) and the scalar dynamic baseline (Linear Router) across batch sizes $B \in \{1, 8, 16, 32, 64\}$ on a single Nvidia GPU. The results are presented in Table~\ref{tab:profiling_results}.

\begin{table*}[t]
\centering
\caption{\textbf{GPU Memory and Inference Latency Profiling.} We report the wall-clock forward pass latency (ms) and peak GPU memory usage (MB) for different batch sizes. Relative memory overhead is computed with respect to the static OFS-Tune baseline.}
\label{tab:profiling_results}
\vskip 0.15in
\begin{tabular}{llccc}
\toprule
\textbf{Batch Size} & \textbf{Model} & \textbf{Latency (ms)} & \textbf{Peak GPU Memory (MB)} & \textbf{Relative Memory Overhead} \\
\midrule
$B=1$ & OFS-Tune (Static) & 8.99 ms & 486.92 MB & - \\
$B=1$ & Linear Router & 11.74 ms & 486.36 MB & $-0.56$ MB ($-0.1\%$) \\
$B=1$ & EpiMerge (Ours) & 17.36 ms & 489.73 MB & $+2.81$ MB ($+0.6\%$) \\
\midrule
$B=8$ & OFS-Tune (Static) & 9.01 ms & 502.34 MB & - \\
$B=8$ & Linear Router & 12.04 ms & 511.46 MB & $+9.12$ MB ($+1.8\%$) \\
$B=8$ & EpiMerge (Ours) & 19.42 ms & 521.11 MB & $+18.77$ MB ($+3.7\%$) \\
\midrule
$B=16$ & OFS-Tune (Static) & 9.09 ms & 520.17 MB & - \\
$B=16$ & Linear Router & 12.13 ms & 538.49 MB & $+18.32$ MB ($+3.5\%$) \\
$B=16$ & EpiMerge (Ours) & 19.47 ms & 556.49 MB & $+36.32$ MB ($+7.0\%$) \\
\midrule
$B=32$ & OFS-Tune (Static) & 9.11 ms & 557.39 MB & - \\
$B=32$ & Linear Router & 12.69 ms & 594.03 MB & $+36.64$ MB ($+6.6\%$) \\
$B=32$ & EpiMerge (Ours) & 19.57 ms & 629.88 MB & $+72.49$ MB ($+13.0\%$) \\
\midrule
$B=64$ & OFS-Tune (Static) & 9.12 ms & 630.78 MB & - \\
$B=64$ & Linear Router & 12.65 ms & 702.83 MB & $+72.05$ MB ($+11.4\%$) \\
$B=64$ & EpiMerge (Ours) & 27.34 ms & 774.83 MB & $+144.05$ MB ($+22.8\%$) \\
\bottomrule
\end{tabular}
\end{table*}

The empirical profiling results confirm our theoretical analysis. While EpiMerge incurs a memory overhead due to the generation and retention of sample-specific weight matrices of shape $[B, D_{out}, D_{in}]$, the absolute overhead remains exceptionally mild. Even at a large deployment batch size of $B=64$, EpiMerge consumes only 774.83 MB of peak GPU memory, which represents an overhead of just 144.05 MB (+22.8\%) compared to the static model. The wall-clock latency also reflects the cost of tensor reconstruction, dynamic contraction, and the forward pass of the frozen sensory extractor copy: EpiMerge takes 27.34 ms at $B=64$ compared to 9.12 ms for the static compromise. This 3x increase in latency and the associated peak memory overhead represent the primary physical trade-offs of true coordinate-wise ensembling.

\subsection{Verification of Test-Time Adaptation and Routing Dynamics}
A critical concern in dynamic model merging is whether the routers and reader heads are actually adapting to the inputs, or if they are simply under-optimized and have converged to a static, task-independent compromise. To verify the presence of active test-time dynamics, we input 20 random samples from each task domain and extract the average learned routing weights of the Linear Router and the epigenetic row/column gating intensities of EpiMerge. The task-specific gating patterns are presented in Table~\ref{tab:routing_analysis}.

\begin{table*}[t]
\centering
\caption{\textbf{Task-Specific Adaptation and Routing Dynamics.} We report the average expert ensembling weights of the Linear Router and the average row-gating intensities of EpiMerge's Epigenetic Reader Heads across different input task distributions.}
\label{tab:routing_analysis}
\vskip 0.15in
\begin{tabular}{lcccc}
\toprule
\textbf{Input Task} & \textbf{MNIST Expert Gating} & \textbf{FashionMNIST Expert Gating} & \textbf{CIFAR-10 Expert Gating} & \textbf{SVHN Expert Gating} \\
\midrule
\multicolumn{5}{l}{\textit{Linear Router Routing Coefficients ($\alpha$)}}\\
MNIST & 0.611 & 0.497 & 0.501 & 0.426 \\
FashionMNIST & 0.557 & 0.499 & 0.500 & 0.463 \\
CIFAR-10 & 0.505 & 0.500 & 0.500 & 0.497 \\
SVHN & 0.505 & 0.500 & 0.499 & 0.498 \\
\midrule
\multicolumn{5}{l}{\textit{EpiMerge Row Gating Intensities ($\mathbf{r}$)}}\\
MNIST & 0.516 & 0.513 & 0.515 & 0.498 \\
FashionMNIST & 0.509 & 0.507 & 0.509 & 0.499 \\
CIFAR-10 & 0.502 & 0.502 & 0.502 & 0.500 \\
SVHN & 0.502 & 0.502 & 0.502 & 0.500 \\
\bottomrule
\end{tabular}
\end{table*}

The gating metrics in Table~\ref{tab:routing_analysis} provide strong empirical evidence of task-sensitive test-time adaptation. 
For the Linear Router, MNIST inputs assign a significantly higher weight to the MNIST expert (0.611) compared to other experts, while SVHN and CIFAR-10 inputs focus their ensembling weights close to 0.50. Interestingly, FashionMNIST inputs heavily activate both the MNIST expert (0.557) and the FashionMNIST expert (0.499), reflecting the high structural and representational similarity between these two grayscale datasets.

For EpiMerge, the Epigenetic Reader Heads exhibit similar task-sensitive behavior. Under MNIST inputs, the row-gating intensity for the MNIST expert reaches 0.516, which is higher than for SVHN (0.498). Under SVHN inputs, the gating intensities converge symmetrically towards 0.502. 
While the gating values remain relatively close to 0.50, this is a natural consequence of the Sigmoid activation and the conservative weight initialization (standard deviation of 0.02) under our short calibration budget. This confirms that test-time routing is indeed dynamic and successfully learns to align parameter scaling with the input task's representational characteristics.

\subsection{Task-Conditioning Oracle and Deployment Limitations}
A major, standard methodology simplification in current multi-task model-merging evaluations (including this study and relevant literature) is the assumption of a \textbf{Task-Conditioning Oracle} at test-time. Specifically, during evaluation, the model loops through the batch and selects the task-specific classification head using the ground-truth task labels:
\begin{equation}
    y_b = \mathcal{H}_{k_{gt}}\left(\mathbf{f}_b\right)
\end{equation}
where $\mathcal{H}_k$ is the classification head for task $k$, and $k_{gt}$ is the ground-truth task index.

In a realistic production deployment scenario, the model does not have access to the ground-truth task label of each test sample. Real-world multi-task deployment requires either a single, unified classification head (spanning all classes of all tasks) or an integrated, non-oracle task-inference module. Assuming perfect oracle guidance at test-time hides the representation overlap and category-wise interference that would occur in a unified class space.

To enhance the transparency and applicability of EpiMerge for production, we propose two concrete architectural pathways to transition to a non-oracle, fully autonomous deployment setup:
\begin{enumerate}
    \item \textbf{Integrated Task Classifier (Two-Stage Inference):} We can train a highly lightweight, frozen task-classification head (e.g., a 2-layer MLP) on top of our frozen sensory extractor's representational vector $\mathbf{h}(x)_b$ using our 64-sample calibration dataset. At test-time, for each sample, the task classifier first predicts the task index:
    \begin{equation}
        \hat{k} = \text{argmax}\left(\text{Softmax}\left(\text{MLP}\left(\mathbf{h}(x)_b\right)\right)\right)
    \end{equation}
    The predicted task index $\hat{k}$ is then used to select the correct active classification head $\mathcal{H}_{\hat{k}}$. Because our sensory extractor is pre-trained and frozen, training this 2-layer task classifier requires negligible parameters and can easily achieve near-perfect task classification on diverse domains.
    \item \textbf{Shared Unified Multi-Task Head:} We can replace the list of task-specific heads with a single, unified classification head $\mathcal{H}_{unified} \in \mathbb{R}^{D_{out} \times C_{total}}$, where $C_{total} = \sum_k C_k$ is the total sum of classes across all merged domains (e.g., 40 categories for our MNIST, FashionMNIST, CIFAR-10, and SVHN mixture). This unified head is then optimized concurrently with the Epigenetic Reader Heads on the offline calibration dataset. This completely bypasses task-conditioned head selection and allows the model to map outputs directly to a single, unified multi-class representation space, handling representation conflicts and class-overlap in a scientifically unified manner.
\end{enumerate}
By outlining these transitions, we bridge the gap between academic task-specific ensembling benchmarks and the practical requirements of end-to-end multi-task AI systems.
